% answers_short.tex —— 2008 年·春季上海卷 · 简短答案速查
% 由 tools/make_short_answers.py 自 answers.tex 的 \gksolhead/\gkanswer 提取，
% 并与 2008/2008_数学真题_春季上海_答案要点.md 逐题比对后生成，勿手改（改源 MD 后重跑）。

\begin{gknotice}
\textbf{使用说明}：本节只给最终答案便于快速核对，完整推导见解析版。答案为编写者独立推导，\textbf{非官方标准答案}。
\end{gknotice}

\gkpart{填空题}
\begin{itemize}[leftmargin=2.6em]
  \item[\textbf{1.}] $(-\infty,4)$
  \item[\textbf{2.}] $\dfrac13$
  \item[\textbf{3.}] $[-2,1)\cup(1,3]$
  \item[\textbf{4.}] $\dfrac{7\pi}{12}$
  \item[\textbf{5.}] $a_n=2n-1$
  \item[\textbf{6.}] $\cos\alpha$
  \item[\textbf{7.}] $5$
  \item[\textbf{8.}] U（按题面文字计算为 $1-\dfrac{\sqrt2}{6}$；固定副本答案页印载 $1+\dfrac{\sqrt2}{6}$）
  \item[\textbf{9.}] $-1$
  \item[\textbf{10.}] $\dfrac1{12}$
  \item[\textbf{11.}] $c_k=a_k-a_{k-1}$
  \item[\textbf{12.}] $\dfrac32$
\end{itemize}

\gkpart{选择题}
\begin{center}\begin{tabular}{ccccc}
\toprule
题号 & 13 & 14 & 15 & 16 \\
\midrule
答案 & C & D & A & B \\
\bottomrule
\end{tabular}\end{center}

\gkpart{解答题}
\begin{itemize}[leftmargin=2.6em]
  \item[\textbf{17.}] $-\dfrac{\sqrt{14}}2$
  \item[\textbf{18.}] $\dfrac{7\sqrt2}{8}$
  \item[\textbf{19.}] (1) $f$ 在实数域严格递增；(2) $m\in[\log_2\dfrac13,\log_2\dfrac35]$
  \item[\textbf{20.}] (1) 上段与下段之比约为 $0.63$；(2) 按 $A,B,C$ 次序，三根钢管长度约为 $60.8$ cm、$36.1$ cm、$60.8$ cm，即 $10\sqrt{37}$ cm、$10\sqrt{13}$ cm、$10\sqrt{37}$ cm
  \item[\textbf{21.}] (1) 是 T 点列；(2) 钝角三角形；(3) $\overrightarrow{A_nA_q}\cdot\vec j>\overrightarrow{A_mA_p}\cdot\vec j$
  \item[\textbf{22.}] (1) $P_z$ 在圆 $C_1:(x-1)^2+y^2=1$ 上；(2) $s$ 为端点 $(-m-r,(m+r)^2)$ 与 $(-m+r,(r-m)^2)$ 间的闭线段；非端点对应非实根；(3) 平行：$m=1$（不同直线时 $r\ne1$）；$s$ 所在直线平分 $s_1$：$r^2=(m-1)^2+1$（$m\ne1$）；等长：$(1+4m^2)r^2=5$
\end{itemize}
